\documentclass[11pt,a4paper]{article}
\usepackage[margin=1in]{geometry}
\usepackage{times}
\usepackage{amsmath,amssymb}
\usepackage{graphicx}
\usepackage{booktabs}
\usepackage{hyperref}
\usepackage[inline]{enumitem}
\usepackage{xcolor}
\usepackage{CJKutf8}

\title{Smart Meeting Assistant: \\ A Bilingual Pipeline with LLM-as-Judge Evaluation}
\author{CS4300 NLP Group Project \\ City University of Hong Kong}
\date{April 2026}

\begin{document}
\maketitle

\begin{abstract}
We present an end-to-end smart meeting assistant for bilingual (Chinese--English) meetings,
unifying five modules --- automatic speech recognition with speaker diarization, summarization,
translation, action-item extraction, and sentiment/engagement analysis --- over a single shared
transcript representation. Each language-understanding module pairs a large-language-model (LLM)
instantiation with a task-specific pretrained baseline (BART, mT5, NLLB-200). We further
introduce a unified LLM-as-Judge evaluation layer and a speaker-constrained function-calling
schema for action-item extraction. Experiments on AMI, AliMeeting, MeetingBank, and MELD show
that the LLM pipeline is competitive with or exceeds specialized baselines on four of five
modules, while the LLM-as-Judge scores correlate with ROUGE and F1 at the meeting level.
\end{abstract}

\section{Introduction}
\label{sec:intro}
% TODO: bilingual meetings are common in HK/academic settings. Five required modules.
% Contributions:
% 1. Unified shared-transcript architecture enabling parallel module execution.
% 2. Bilingual-aware prompting with speaker context window.
% 3. LLM-as-Judge evaluation across summary / actions.
% 4. Comparative pretrained-vs-LLM baselines for 4 of 5 modules.
% 5. A Gradio demo with retrieval-augmented Q&A.

\section{Related Work}
\label{sec:related}
% TODO: Whisper (Radford et al.); pyannote (Bredin et al.); AMI-based summarization
% (Shang et al.; Zhang et al.); MeetingBank (Hu et al.); NLLB (NLLB team 2022);
% LLM-as-Judge (Zheng et al. 2023); Self-RAG (Asai et al. 2024).

\section{Methodology}
\label{sec:method}

\subsection{Shared-Transcript Architecture}
% TODO: architecture figure (fig:arch) showing ASR -> transcript JSON -> 4 parallel modules + QA.

\subsection{Bilingual-Aware Prompting}
% TODO: describe language auto-detection, target-lang flag, speaker-aware context window for translation.

\subsection{Structured Outputs via Function Calling}
% TODO: action-item JSON schema; assignee enum constrained to speakers produced by Mod 1.

\subsection{LLM-as-Judge}
% TODO: 1-5 Likert on faithfulness / coverage / specificity; correlation to ROUGE / F1.

\section{Experiments}
\label{sec:exp}

\subsection{Datasets}
\begin{table}[h]
\centering
\begin{tabular}{lllr}
\toprule
Dataset & Language & Purpose & Size \\
\midrule
AMI Meeting Corpus & EN & M1--5 ground truth & 10 meetings \\
AliMeeting (M2MeT) & ZH & M1, M3 & 5 sessions \\
MeetingBank & EN & M2 (ROUGE) & 50 summaries \\
MELD & EN & M5 (emotion) & 200 utts \\
\bottomrule
\end{tabular}
\caption{Evaluation datasets.}
\end{table}

\subsection{Module 1: ASR + Diarization}
% TODO: WER (EN), CER (ZH), DER. Table.

\subsection{Module 2: Summarization}
% TODO: ROUGE-1/2/L, BERTScore. Table comparing GPT-4o-mini vs BART vs mT5.

\subsection{Module 3: Translation}
% TODO: sacreBLEU. Table: GPT-4o-mini vs NLLB-200 on ZH->EN and EN->ZH.

\subsection{Module 4: Action Items}
% TODO: P/R/F1 + assignee accuracy on 30 AMI segments; Cohen's kappa for annotator agreement.

\subsection{Module 5: Sentiment}
% TODO: MELD accuracy / macro-F1. Show timeline figure on AMI ES2004a.

\subsection{LLM-as-Judge}
% TODO: Table of Likert averages per system per kind. Correlation plot with automated metrics.

\section{Discussion}
\label{sec:disc}
% TODO: Where LLM wins (summary novelty, action items). Where baselines are competitive (MT BLEU).
% Bilingual code-switching failures. Cost and latency analysis (cache helps).

\section{Conclusion}
\label{sec:conclusion}
% TODO: Wrap up. Limitations + future work (on-device ASR, long-context summarization, humans-in-the-loop).

\bibliographystyle{plain}
\bibliography{refs}

\appendix
\section{Prompts}
% TODO: paste full system prompts used for each module.

\section{Presentation slides}
% TODO: \includepdf[pages=-]{../slides/slides.pdf}

\end{document}
